\documentclass[11pt,a4paper]{article}
\usepackage[T1]{fontenc}
\usepackage[utf8]{inputenc}
\usepackage{lmodern,microtype}
\usepackage[margin=23mm]{geometry}
\usepackage{array,booktabs,longtable,tabularx}
\usepackage[table]{xcolor}
\usepackage{enumitem}
\usepackage[authoryear,round]{natbib}
\usepackage{hyperref}
\hypersetup{colorlinks=true,linkcolor=blue!45!black,citecolor=blue!45!black,urlcolor=blue!45!black,pdftitle={Requirements Specification for a Knowledge Graph Exploration Application}}
\urlstyle{same}
\setlength{\parindent}{0pt}
\setlength{\parskip}{5pt}
\setlength{\tabcolsep}{4pt}
\renewcommand{\arraystretch}{1.16}
\setlist[itemize]{nosep,leftmargin=*}
\setlist[enumerate]{nosep,leftmargin=*}
\newcolumntype{P}[1]{>{\raggedright\arraybackslash}p{#1}}
\newcommand{\req}[4]{\subsubsection*{#1 --- #2}\phantomsection\label{req:#1}\textbf{Priority:} #3.\quad\textbf{Tests:} #4.\par}
\newcommand{\origin}[3]{\par{\small\textbf{Sources (#1).} #2 \textbf{Derivation.} #3\par}}
\newcommand{\evidence}[1]{\hyperlink{evidence:#1}{#1}}
\title{Requirements Specification for a\\Knowledge Graph Exploration Application}
\author{}
\date{Version 1.4\\6 October 2026\\Literature cutoff: 29 August 2026}
\begin{document}
\maketitle
\begin{abstract}
This specification defines a domain-independent application for researchers and users without graph-query expertise. Version 1 supports both Resource Description Framework (RDF) graphs queried through SPARQL and property graphs queried through Cypher. Users search for entities, refine results with facets, navigate relationships, find paths, construct queries with guidance, switch views, and save findings with their source and query context. Requirements are linked to primary literature and paired acceptance scenarios. Numerical thresholds are proposed engineering targets, not measured results for the application.
\end{abstract}
\tableofcontents
\newpage

\section{Purpose, scope, and success criteria}
The application shall let users revise search terms, filters, and relationship constraints while inspecting graph data. The core tasks are to find entities, inspect connections, restrict a result set, formulate a query, and save a finding with its supporting data. Researchers shall inspect native queries and source values. Non-experts shall complete the core workflow without writing SPARQL or Cypher.

Version 1 requires RDF/SPARQL and property-graph/Cypher connectors for every mandatory exploration workflow and the visual query profile in Section~\ref{sec:profile}.

The application shall operate independently of a domain vocabulary. Source profiling may inspect labels, types, relationships, and value types. Users may configure presentation labels and preferred facets without changing application code. The application explores external graph sources.

Source connections are read-only. Version 1 excludes graph construction and editing, ontology alignment, bulk ingestion into an application-owned store, autonomous external actions, cross-source query federation, arbitrary language conversion, and persistent collaboration services. Users may edit workspace annotations and saved queries locally.

\subsection{Priority conventions}
\textbf{Must} requirements are release conditions. \textbf{Should} requirements improve the release but cannot compensate for a missing Must capability. Functional requirements (FR), non-functional requirements (NFR), and acceptance scenarios (AC) use stable identifiers; a changed requirement shall receive a new version and updated tests. The requirements specify behaviour; implementation technology is left open.

\subsection{Document status and delivery boundary}
Version 1.4 is the completed behavioural specification for implementation and acceptance planning. It specifies required behaviour and proposed targets; it does not certify that an application has been built or tested. The survey companion website is a separate deliverable for reading and comparing research reports. Its catalogue filters are not a graph-query builder, and publication of that website does not satisfy application acceptance.

All Must requirements belong to the version 1 baseline. Implementation may proceed in increments: connections and semantic fidelity; search, facets and navigation; visual/native queries and inspectable guidance; coordinated views, workspace recovery and exports; then paired acceptance, accessibility and performance evaluation. Intermediate increments are development milestones, not a reduced-compliance version 1 release. The specification introduces no delivery date or budget commitment.

\section{Evidence basis and design rationale}
The literature cutoff is 29 August 2026. Table~\ref{tab:evidence} lists eight research papers, four technical specifications, and one security-guidance source, with section references and source URLs. The papers inform the proposed features. Mandatory support for both graph models, priorities, and numerical thresholds are project decisions.

Each origin statement distinguishes \emph{documented precedent} (behaviour described in a paper), \emph{normative basis} (adopted language or accessibility semantics), \emph{design extension} (features added to the precedent), and \emph{project decision} (scope, policy, formats, targets, or tests). A documented feature alone does not establish a usability benefit. Citations support the stated precedent or semantics. Evidence identifiers link to the source register, whose URLs open the cited paper or technical document. Acceptance scenarios identify the same underlying sources; their fixtures and pass conditions are application-specific tests.

Tabulator combines exploration and analysis with source inspection \citep[Sections 3.1--3.4]{tabulator}. /facet describes dynamic facets, search, and cross-type constraints \citep[Sections 3--4]{facet}; RelFinder addresses disambiguation and relationship search \citep[Sections 2.1--2.3]{relfinder}. Sparklis combines stepwise query refinement with controlled-language presentation \citep[Sections 3 and 5]{sparklis}.

LODmilla saves graph state and reverses exploration actions \citep[Section 3.1]{lodmilla}. Sampo-UI connects a faceted result set with analytical perspectives \citep[Sections 3.1--3.2 and 4.4--4.5]{sampoui}. Ranked graph-query expansions motivate inspectable guidance \citep[Sections 3--4]{suggestions}. SemSpect explicitly addresses RDF and labelled property graphs and uses aggregation to reduce visual clutter \citep[Abstract and Sections 1--2]{semspect}.

{\small
\begin{longtable}{@{}P{1.1cm}P{14.6cm}@{}}
\caption{Primary evidence register. Locators refer to named or numbered source sections.}\label{tab:evidence}\\
\toprule ID & Primary source, locator, and relevant observation\\ \midrule
\endfirsthead\toprule ID & Primary source, locator, and relevant observation\\ \midrule\endhead
\bottomrule\endlastfoot
\hypertarget{evidence:E01}{E01} & Tabulator \citep{tabulator}, Sections 3.1--3.4: source provenance and network status; exploration and analysis; outliner; query-by-example and alternative result views.\par\smallskip\url{https://www.cs.columbia.edu/~chilton/web/my_publications/BernersLeeTabulator2006.pdf}\\
\hypertarget{evidence:E02}{E02} & /facet \citep{facet}, Sections 3.1--3.5 and 4.1--4.5: dynamic facets, semantic keyword search, multi-type queries, runtime facet specification, and facet-specific displays.\par\smallskip\url{https://doi.org/10.1007/11926078_20}\\
\hypertarget{evidence:E03}{E03} & RelFinder \citep{relfinder}, Sections 2.1--2.3: entity disambiguation, iterative relationship search, labelled directed edges, and graph display control.\par\smallskip\url{https://doi.org/10.1007/978-3-642-10543-2_21}\\
\hypertarget{evidence:E04}{E04} & Sparklis \citep{sparklis}, Section 3, \emph{Principles and architecture}; Section 5, \emph{Capabilities and limitations}; Section 5.5, \emph{Portability: No required configuration}: guided refinement, readable queries, and endpoint-driven vocabulary discovery.\par\smallskip\url{https://doi.org/10.3233/SW-150208}\\
\hypertarget{evidence:E05}{E05} & LODmilla \citep{lodmilla}, Section 3.1, \emph{Frontend}, subheadings \emph{Saving/Sharing/Loading} and \emph{Undo}; Section 3.2, \emph{Backend}: preserved graph state, undo, source traversal, and bounded search.\par\smallskip\url{https://doi.org/10.1007/978-3-319-08425-1_9}\\
\hypertarget{evidence:E06}{E06} & Sampo-UI \citep{sampoui}, Sections 3.1--3.2 and 4.1, 4.4--4.5: facets, shared result perspectives, interface structure, and CSV query-result export.\par\smallskip\url{https://doi.org/10.3233/SW-210428}\\
\hypertarget{evidence:E07}{E07} & Graph-query suggestions \citep{suggestions}, Section 3, \emph{Problem Formulation}; Section 4, \emph{Graph-Query Suggestion Model}: exemplar queries and ranked edge expansions.\par\smallskip\url{https://doi.org/10.1145/3366423.3380005}\\
\hypertarget{evidence:E08}{E08} & SemSpect \citep{semspect}, Abstract; Section 1, \emph{Motivation}; Section 2, \emph{Querying and Visual Exploration Beyond the Hairball}: both graph families, aggregation, and overview-to-detail exploration.\par\smallskip\url{https://ceur-ws.org/Vol-3526/paper-07.pdf}\\
\hypertarget{evidence:E09}{E09} & RDF 1.1 Concepts \citep{rdfconcepts}, Sections 3.1--3.4 and 4: term identity, literal structure, blank-node scope, and datasets.\par\smallskip\url{https://www.w3.org/TR/2014/REC-rdf11-concepts-20140225/}\\
\hypertarget{evidence:E10}{E10} & SPARQL 1.1 Query Language \citep{sparql}, Sections 5--9 and 11--17: patterns, optionality, alternatives, paths, aggregates, result modifiers, and value operations. The selected visual subset is a project decision.\par\smallskip\url{https://www.w3.org/TR/2013/REC-sparql11-query-20130321/}\\
\hypertarget{evidence:E11}{E11} & openCypher \citep{opencypher}, pinned release 2024.3 grammar and TCK: native-language semantics and scenarios with graphs, parameters, expected results or errors.\par\smallskip\url{https://github.com/opencypher/openCypher/releases/tag/2024.3}\\
\hypertarget{evidence:E12}{E12} & WCAG 2.2 \citep{wcag}, Success Criteria 1.1.1, 2.1.1, 2.4.7, 4.1.2 and Section 5.2.1: text alternatives, keyboard access, visible focus, control semantics, and conformance levels.\par\smallskip\url{https://www.w3.org/TR/2024/REC-WCAG22-20241212/}\\
\hypertarget{evidence:E13}{E13} & OWASP REST Security \citep{owasprest}, \emph{Sensitive information in HTTP requests}: avoiding credential exposure in URLs and server logs. This guidance supports part of NFR-06; application storage and export policies remain project decisions.\par\smallskip\url{https://cheatsheetseries.owasp.org/cheatsheets/REST_Security_Cheat_Sheet.html#sensitive-information-in-http-requests}\\
\end{longtable}}

\begin{table}[htbp]\centering\small
\caption{Design considerations and sources.}\label{tab:gaps}
\begin{tabularx}{\textwidth}{@{}P{1cm}P{2.1cm}X@{}}
\toprule ID & Sources & Design consideration\\\midrule
G01 & \evidence{E01}, \evidence{E02}, \evidence{E04} & Vocabulary and query-language knowledge can obstruct entry. Provide readable search, facets, and progressive query construction.\\
G02 & \evidence{E03}, \evidence{E05}, \evidence{E08} & Expansion increases visual density and can lose context. Bound displays, expose paths, and preserve recovery.\\
G03 & \evidence{E04}, \evidence{E07} & Guidance changes exploration direction. Explain its basis and distinguish verified results from estimates.\\
G04 & \evidence{E01}, \evidence{E05} & Findings require source, query, and session context to be revisited. Preserve context in saves and exports.\\
G05 & \evidence{E06}, \evidence{E08} & Set inspection and aggregation require coordinated views beyond a node-link display.\\
G06 & \evidence{E01}, \evidence{E04}, \evidence{E08} & Source access, graph semantics, and limits differ. Test both connectors and their failure modes separately.\\
\bottomrule\end{tabularx}\end{table}

\section{Users and tasks}
\textbf{Researchers} explore unfamiliar data, formulate and inspect queries, compare refinements, check values and scope, and export reproducible evidence. \textbf{Non-experts} start from a phrase or entity, disambiguate labels, follow relationships, refine sets, and preserve findings through readable controls.

Either user may select a configured connection or enter an authorised source profile. A profile contains graph family, address, dataset or graph scope, authentication reference, and query dialect. It does not imply permission to modify the source.

The core sequence is: select source; search or choose a type; inspect entities; refine facets; expand relationships or find a path; build or inspect a query; change view; save a finding; export its supporting context. Users may repeat or branch these steps without completing a fixed wizard.

\subsection{Reference use cases and domain fixtures}
The following are proposed acceptance fixtures, not observations about commissioned users or available institutional datasets. Each fixture shall be represented independently in RDF and a property graph, with model-specific identity and an explicit result oracle.
\begin{description}[leftmargin=0pt,labelindent=0pt,labelwidth=0pt,style=nextline]
\item[UC-01: Research-publication exploration.] A researcher starts with an ambiguous author name, selects the intended person, filters publications by year and topic, follows affiliations or co-authorship, and asks for a bounded relationship between two authors. The researcher inspects the native query, saves a finding and exports its evidence bundle. Distinct people with the same label remain distinct. Jointly authored records do not imply a relationship that is absent from the source. Verify FR-01--10 and FR-13--16 through AC-01--05, AC-07 and AC-08.
\item[UC-02: Cultural-heritage exploration.] A non-expert starts from a person or place label, restricts a collection by date and category, follows a documented relationship, and changes between list and graph views without losing filters. The fixture includes multilingual labels, missing dates, repeated relationship types, and a disconnected entity. The user explains the active constraints and saves a finding without writing a native query. Verify FR-01--07 and FR-09--16 through AC-01--08.
\item[UC-03: Recovery and portability.] Run both information needs with keyboard controls, on both graph families, on a narrow viewport, and after injected connection and authentication failures. Reopen an exported workspace against an updated fixture and identify stale observations before refreshing. Verify NFR-01--11 through AC-07--11. The two domain vocabularies provide the portability test; equivalent content does not require identical storage encodings.
\end{description}

The small correctness fixtures are distinct from the 100,000-entity performance fixture. Every fixture package shall contain source files or reproducible generation instructions, a version/hash, graph scope, connection configuration without secrets, input states, expected results and deliberately failing cases. Personal names and events in synthetic fixtures shall be identified as synthetic. No real participant data are required for these technical tests.

\subsection{Deployment, connection administration, and local data}
The version 1 deployment baseline is a browser interface with explicitly configured read-only graph adapters. It may run on localhost or an institutional HTTPS service. A deployment declares its supported browser/runtime and adapter versions. Public endpoints may be accessed directly when their access policy permits it; an authenticated source or one requiring mediation uses a configured adapter service. Deployment administrators configure allowed source addresses, scopes, dialects and server-side secret references. Users choose or enter authorised profiles within those controls. An application-side read-only parser does not replace source-side read permissions.

Workspace data remain local to the user's browser in version 1, with explicit export/import and deletion. Saving locally is a deliberate user action; the interface states that local storage is specific to the browser profile and may be cleared. Persistent shared workspaces, background telemetry and central participant registries are outside this baseline. Authentication material stays outside saved profiles, workspace files, logs and exports. Reopening a restricted source may require fresh authentication. A user may inspect, download or delete local workspace data without contacting the source.

\section{Graph models, connections, and data semantics}
RDF terms and dataset context follow the RDF 1.1 data model \citep{rdfconcepts}. Resources retain their Internationalized Resource Identifier (IRI) or source-scoped blank-node identity. Literals retain lexical value, datatype, and language tag. Named-graph context remains inspectable when provided. Blank-node labels shall not be treated as persistent cross-session identifiers.

Property-graph nodes retain source-scoped identifiers and all labels. Relationships retain identifiers where supplied, type, direction, endpoints, and properties. Parallel relationships and self-loops remain distinguishable, including when grouped for display.

The shared interface shall not erase model-specific semantics. Missing properties, unbound variables, empty strings, zero, and false remain distinct. Observed schema information is labelled as observed; declared schema retains its source. Source-side read permissions and application operation controls enforce read-only access. SPARQL Update, mutating Cypher statements, and administrative or mutating procedures are rejected.

\subsection{Version 1 visual query profile}\label{sec:profile}
The builder supports connected patterns, variables, repeated-variable joins, type or label constraints, directed and inverse relationships, optional patterns, alternatives, equality and inequality, string matching, numeric and temporal ranges, projection, distinct results, ordering, limits, and grouped \texttt{COUNT}, \texttt{SUM}, \texttt{AVG}, \texttt{MIN}, and \texttt{MAX}. Filters and aggregates check applicable value types.

The RDF connector generates SPARQL 1.1 read queries \citep{sparql}. The property-graph connector generates the supported operations in its declared Cypher dialect. Each connector declares its language and version and retains its model-specific semantics. Native read queries may use additional supported constructs. Queries outside the visual profile remain in native mode; unsupported constructs are not discarded or assigned an invented visual translation. Changing graph family requires a fresh query or explicit supported reconstruction.

\subsection{Dialect compatibility declaration}
Each Cypher connector profile shall identify its implementation dialect and version, its language-reference release and commit, and semantic differences relevant to the mandatory visual query profile. The openCypher specification, grammar, and Technology Compatibility Kit (TCK) provide a versioned reference \citep{opencypher}. The inspected reference is release 2024.3, published on 20 March 2026, at commit \texttt{677cbafabb8c3c5eed458fd3b1ec0daec8d67d23}.

The compatibility manifest shall map every mandatory visual construct to supported native syntax, expected semantics, and its acceptance cases. A connector may implement another declared Cypher dialect; every Must workflow and operator must still pass separately on both graph families. Unsupported mandatory constructs prevent release. Relevant read-query TCK scenarios may inform language checks, but a selected subset does not establish complete TCK conformance. This declaration does not prescribe a database product or require arbitrary SPARQL-to-Cypher translation.

\section{Functional requirements}
\subsection{Connections and graph fidelity}
\par\noindent\begin{minipage}{\linewidth}
\req{FR-01}{Read-only dual-backend connections}{Must}{AC-01, AC-08}
Connect to an RDF/SPARQL endpoint and a property-graph/Cypher source through profiles. Expose source scope, dialect, connection, and authentication status. Changing source preserves the earlier workspace as a separate context. Both connectors implement every mandatory workflow.
\origin{Documented precedent; project decision}{Sparklis, \evidence{E04}, Section 5.5 \citep{sparklis}, describes SPARQL endpoint access; SemSpect, \evidence{E08}, Section 2 \citep{semspect}, supports RDF and labelled property graphs.}{Mandatory coverage of both models in version 1, connection profiles, read-only access, authentication feedback, and separate source workspaces are project decisions. SemSpect's indexed RDF dumps do not establish remote endpoint support.}
\end{minipage}\par
\par\noindent\begin{minipage}{\linewidth}
\req{FR-02}{Identity and value preservation}{Must}{AC-01, AC-07}
Details, table cells, and exports preserve the semantics above. Equal labels never merge unequal entities. Identifiers remain inspectable. Language preference changes presentation, not source values.
\origin{Normative basis; design extension}{RDF Concepts, \evidence{E09}, Sections 3.1--3.4 and 4 \citep{rdfconcepts}, defines terms and dataset identity; the declared Cypher reference, \evidence{E11} \citep{opencypher}, governs property-graph semantics.}{Preserving those distinctions in every view and export, and preventing label-based entity merging, are interface obligations derived from the adopted semantics. Language preference remains a presentation choice.}
\end{minipage}\par

\subsection{Search, schema inspection, and facets}
\par\noindent\begin{minipage}{\linewidth}
\req{FR-03}{Entity search and disambiguation}{Must}{AC-02}
Search by identifier or label text. Candidates expose type or label, source, and identifier. State search scope and result limits. Expose discovered entity types or labels, relationships, and property value types.
\origin{Documented precedent; design extension}{/facet, \evidence{E02}, Section 4.2 \citep{facet}, describes keyword search; RelFinder, \evidence{E03}, Section 2.1 \citep{relfinder}, describes entity disambiguation; Sparklis, \evidence{E04}, Section 5.5 \citep{sparklis}, discovers endpoint vocabulary.}{The application adds candidate identifiers, source and type information, search bounds, and property-value inspection.}
\end{minipage}\par
\par\noindent\begin{minipage}{\linewidth}
\req{FR-04}{Faceted refinement}{Must}{AC-02}
Combine categorical, numeric, and temporal facets. Selections within a categorical facet use OR; selections across facets use AND. Show removable constraints and exact or estimated counts. Empty results preserve constraints and support removal or undo. Missing values remain distinct from empty strings.
\origin{Documented precedent; design extension; project decision}{/facet, \evidence{E02}, Sections 3--4 \citep{facet}, supplies dynamic facets; Sampo-UI, \evidence{E06}, Section 4.4 \citep{sampoui}, supplies active filters, disjunctive choices and date/integer ranges.}{OR within and AND across facets is the adopted application contract. Estimate labels, empty-result recovery and missing-value distinctions specify how that contract is exposed.}
\end{minipage}\par

\subsection{Navigation and relationship discovery}
\par\noindent\begin{minipage}{\linewidth}
\req{FR-05}{Incremental neighbourhood navigation}{Must}{AC-03}
Inspect incoming and outgoing relationships; expand a relationship type; pin an entity; collapse and restore expansions. Expansion is bounded and pageable. Collapse hides view elements while preserving source data and exploration history. Edges retain direction and type.
\origin{Documented precedent; design extension}{Tabulator, \evidence{E01}, Section 3.3 \citep{tabulator}, exposes incoming/outgoing links; RelFinder, \evidence{E03}, Section 2.3 \citep{relfinder}, supports directed labelled graph inspection and pinning; LODmilla, \evidence{E05}, Section 3.1 \citep{lodmilla}, opens related resources and removes canvas nodes.}{Pageable expansion, restoration after collapse, and preservation of source and history state define the proposed workspace behaviour.}
\end{minipage}\par
\par\noindent\begin{minipage}{\linewidth}
\req{FR-06}{Bounded relationship paths}{Must}{AC-03}
Search between two selected entities with relationship and direction constraints. Default depth is three hops, configurable up to six. Initially show at most ten paths ordered by hop count and a deterministic tie-break. Expose path identities, directions, and source scope. A reached limit is not proof that no path exists.
\origin{Documented precedent; design extension; project decision}{RelFinder, \evidence{E03}, Section 2.2 \citep{relfinder}, searches paths with configurable maximum length; LODmilla, \evidence{E05}, Section 3.2 \citep{lodmilla}, bounds traversal and results.}{Three default hops, six maximum hops, ten initial paths and deterministic ordering are selected project limits. The application adds direction controls and messages that distinguish a search limit from a missing path.}
\end{minipage}\par

\subsection{Visual and native query construction}
\par\noindent\begin{minipage}{\linewidth}
\req{FR-07}{Visual query profile}{Must}{AC-04}
Implement Section~\ref{sec:profile} for both models. Add, change, and remove constraints incrementally. Show generated query text, parameters, scope, and a readable description before execution.
\origin{Documented precedent; normative basis; project decision}{Tabulator, \evidence{E01}, Section 3.4 \citep{tabulator}, describes query formulation; Sparklis, \evidence{E04}, Sections 3 and 5.1 \citep{sparklis}, supports incremental guided construction. \evidence{E10} and \evidence{E11} provide native-language semantics \citep{sparql,opencypher}.}{The exact operator profile and mandatory coverage on both models are project decisions. Showing query text, parameters, scope and explanation before execution is the proposed interface contract.}
\end{minipage}\par
\par\noindent\begin{minipage}{\linewidth}
\req{FR-08}{Expert read-query mode}{Must}{AC-04, AC-08}
Inspect, edit, execute, and save native read queries with language and dialect. Native editing does not overwrite the saved visual state silently. Errors identify rejected syntax or capabilities. Queries outside the visual profile remain usable in native mode.
\origin{Documented precedent; design extension; project decision}{Tabulator, \evidence{E01}, Section 3.4 \citep{tabulator}, permits viewing and editing SPARQL query text.}{Separate native and visual state, saved native queries, capability-specific errors and retention of constructs outside the visual subset extend that precedent. Native Cypher mode and read-only enforcement are scope decisions; the paper does not establish arbitrary cross-language translation.}
\end{minipage}\par

\subsection{Explainable guidance}
\par\noindent\begin{minipage}{\linewidth}
\req{FR-09}{Inspectable next-step suggestions}{Must}{AC-05}
Provide user-requested eligible relationships, refinements, or result views. Each suggestion states its basis: schema compatibility, checked counts, or ranking criteria. Distinguish schema-valid, estimated, and verified non-empty suggestions. Users may ignore guidance.
\origin{Documented precedent; design extension}{Sparklis, \evidence{E04}, Sections 3 and 5.2 \citep{sparklis}, supplies data-driven refinements; Graph-Query Suggestions, \evidence{E07}, Sections 3--4 \citep{suggestions}, ranks graph-query expansions.}{The application explains each suggestion's basis and distinguishes schema compatibility, estimates and verified non-empty results. Result-view suggestions extend the cited query-expansion mechanisms.}
\end{minipage}\par
\par\noindent\begin{minipage}{\linewidth}
\req{FR-10}{Readable query descriptions}{Must}{AC-05}
Describe active query and filters in plain English from the current query state. Distinguish conjunction, alternatives, optional values, aggregates, and limits. Descriptions use source labels and do not imply stronger findings than the query produces.
\origin{Documented precedent; design extension; project decision}{Sparklis, \evidence{E04}, Section 5.3 \citep{sparklis}, verbalizes query state with source labels in English or French.}{Selecting English is a project decision. Faithful descriptions of every selected operator, its scope and limits are correctness obligations derived for this application.}
\end{minipage}\par

\subsection{Coordinated views and complexity control}
\par\noindent\begin{minipage}{\linewidth}
\req{FR-11}{Graph, table, and detail coordination}{Must}{AC-06}
Views use the same result scope and preserve selection, filters, and query context. A selected entity is identifiable across views. Textual controls provide equivalent core inspection and navigation when the graph cannot be used.
\origin{Documented precedent; normative basis; design extension}{Tabulator, \evidence{E01}, Section 3.4 \citep{tabulator}, presents alternative query-result views; Sampo-UI, \evidence{E06}, Sections 4.4--4.5 \citep{sampoui}, coordinates facet and result state. \evidence{E12} grounds accessible interaction \citep{wcag}.}{Preserving entity selection across views and offering equivalent textual operation are application-level coordination and accessibility adaptations.}
\end{minipage}\par
\par\noindent\begin{minipage}{\linewidth}
\req{FR-12}{Aggregation and display limits}{Must}{AC-06, AC-09}
Show counts and grouped summaries with computation scope and completeness. Drill from an aggregate to its entities. Initially display at most 500 nodes and 1,000 relationships; use grouping, paging, or selected slices beyond this envelope. Display omission never changes the reported query result.
\origin{Documented precedent; design extension; project decision}{SemSpect, \evidence{E08}, Sections 1--2 \citep{semspect}, uses aggregation and overview/detail; Sampo-UI, \evidence{E06}, Sections 3.1 and 4.5 \citep{sampoui}, provides summaries, paging and clustering.}{The display limits of 500 nodes and 1,000 relationships are project targets. Completeness labels, aggregation scope and query totals unchanged by display limits are application extensions.}
\end{minipage}\par
\par\noindent\begin{minipage}{\linewidth}
\textbf{FR-12a (Should).} Offer temporal and geographical views when source values support them without inventing dates, coordinates, or units. Otherwise retain table inspection.
\origin{Documented precedent; design extension; project decision}{Tabulator, \evidence{E01}, Section 3.4 \citep{tabulator}, includes map/calendar views; /facet, \evidence{E02}, Section 4.5 \citep{facet}, includes a timeline; Sampo-UI, \evidence{E06}, Section 4.5 \citep{sampoui}, includes geographic and temporal perspectives.}{Should priority, value-availability checks, preservation of units and the table fallback are choices for this domain-independent application.}
\end{minipage}\par

\subsection{Workspace, provenance, and export}
\par\noindent\begin{minipage}{\linewidth}
\req{FR-13}{Recoverable exploration workspace}{Must}{AC-07}
Save and reopen queries, filters, selections, annotations, positions, source references, and steps. Undo and redo restore view and query state. Branching preserves earlier branches. Reopening distinguishes saved observations from refreshed results.
\origin{Documented precedent; design extension}{LODmilla, \evidence{E05}, Section 3.1, Saving/\allowbreak Sharing/\allowbreak Loading and Undo \citep{lodmilla}, preserves graph state and reverses an action; Sparklis, \evidence{E04}, Section 4 \citep{sparklis}, provides query permalinks.}{Redo, branching, annotations, additional workspace fields, and the distinction between saved and refreshed results extend the cited capabilities.}
\end{minipage}\par
\par\noindent\begin{minipage}{\linewidth}
\req{FR-14}{Data and query export}{Must}{AC-07}
Export tables as comma-separated values (CSV) and typed JavaScript Object Notation (JSON), native query text, and graph views as Scalable Vector Graphics (SVG). RDF statement exports retain graph context; property-graph exports retain identities and properties. Indicate full versus bounded results. Typed JSON preserves distinctions flattened for CSV presentation.
\origin{Documented precedent; normative basis; project decision}{Sampo-UI, \evidence{E06}, Section 4.1 \citep{sampoui}, exports CSV query results; SemSpect, \evidence{E08}, Section 2 \citep{semspect}, exports SPARQL/Cypher queries. \evidence{E09} and \evidence{E11} ground model fidelity \citep{rdfconcepts,opencypher}.}{CSV, typed JSON and SVG are selected output formats. Preserving typed identities and graph context, and disclosing bounds, are derived export obligations.}
\end{minipage}\par
\par\noindent\begin{minipage}{\linewidth}
\req{FR-15}{Finding provenance}{Must}{AC-07}
An evidence bundle includes finding, source profile identifier, graph scope, retrieval time, native query and parameters, filters, limits, available source version, and annotations. An unavailable version is recorded as unknown. Provenance does not imply independent verification of returned assertions.
\origin{Documented precedent; design extension}{Tabulator, \evidence{E01}, Section 3.1 \citep{tabulator}, exposes source provenance; LODmilla, \evidence{E05}, Section 3.1 \citep{lodmilla}, saves graph state representing findings.}{The source, query, time and version fields in the evidence bundle extend the cited features.}
\end{minipage}\par

\subsection{Recovery and evaluation support}
\par\noindent\begin{minipage}{\linewidth}
\req{FR-16}{Cancellation and recoverable failure}{Must}{AC-08}
Cancel or retry requests without losing the last valid workspace. Distinguish connection, authentication, query, timeout, and decoding failures. Late cancelled results do not replace current state. Retained results show timestamp and stale status.
\origin{Documented precedent; design extension}{Tabulator, \evidence{E01}, Section 3.1 \citep{tabulator}, exposes network progress and failures; Sparklis, \evidence{E04}, Sections 5.4--5.5 \citep{sparklis}, describes endpoint-dependent performance and limits.}{Cancellation, retry, differentiated failures, retained stale results and suppression of late responses define this application's recovery contract.}
\end{minipage}\par
\par\noindent\begin{minipage}{\linewidth}
\textbf{FR-17 (Should).} Optional evaluation recording captures task identifiers, enabled interaction events, completion, time, and errors. It excludes credentials and unnecessary data values. Users inspect, export, and delete records. Recording is disabled by default.
\origin{Documented precedent; project decision}{Sparklis, \evidence{E04}, Section 6 introductory paragraphs \citep{sparklis}, describes navigation logs and a user logging switch.}{Task/error instrumentation, inspect/export/delete controls, data minimization and recording disabled by default are project decisions. The precedent's IP/query/endpoint logging does not supply these privacy defaults.}
\end{minipage}\par

\subsection{Source profiles and workspace interchange}
\par\noindent\begin{minipage}{\linewidth}
\req{FR-18}{Inspectable source configuration}{Must}{AC-10}
Before exploration, show the selected graph family, source profile identifier, graph scope, declared query dialect/version and connection state. Distinguish configured labels and facets from source-observed schema. Validate a profile before saving it; report unsupported mandatory operations before query construction. Changing the profile does not silently execute an old query against a new source. Credentials are referenced separately and remain outside the saved profile.
\origin{Design extension; project decision}{Sparklis, \evidence{E04}, Section 5.5 \citep{sparklis}, describes endpoint-driven discovery; the model and language semantics are grounded in \evidence{E09}--\evidence{E11}.}{The visible configuration contract, preflight, capability declaration and source-change behaviour are project decisions that make FR-01 and FR-07 testable.}
\end{minipage}\par
\par\noindent\begin{minipage}{\linewidth}
\req{FR-19}{Versioned workspace import and deletion}{Must}{AC-11}
Export and import a versioned workspace file containing the fields in FR-13 and FR-15 and the relevant data-model identifier. Validate schema version, identifiers, field types and size before importing. Preview source/scope differences; require an explicit choice to create a separate workspace or replace the current one. Invalid or unsupported files leave the current workspace intact. Imported query text is inert until an explicit execution action; annotations are displayed as text. Allow deletion of a saved workspace and associated local evaluation records, with a clear indication of what will be removed. Credential material is never restored from an imported file.
\origin{Design extension; project decision}{LODmilla, \evidence{E05}, Section 3.1 \citep{lodmilla}, motivates recoverable saved state; Tabulator, \evidence{E01}, Section 3.1 \citep{tabulator}, motivates retained source context.}{The versioned interchange schema, validation, explicit replacement, inert imports and local deletion contract extend these precedents.}
\end{minipage}\par

\section{Quality requirements and operational constraints}
The timings below are engineering targets under the declared benchmark conditions. Measure local user interface (UI) feedback separately from source-query execution.

{\small
\begin{longtable}{@{}P{1.35cm}P{14.35cm}@{}}
\caption{Mandatory quality requirements.}\label{tab:nfr}\\
\toprule ID & Requirement and test\\\midrule
\endfirsthead\toprule ID & Requirement and test\\\midrule\endhead
\bottomrule\endlastfoot
NFR-01 & Exact queries match the fixture oracle for every mandatory case. Preserve identity, direction, datatype, named scope, and aggregate scope. Test both backend families separately.\par\smallskip \textbf{Sources: normative basis and project verification.} \evidence{E09}--\evidence{E11} \citep{rdfconcepts,sparql,opencypher} define the data and query semantics. Fixture oracles and separate paired acceptance are the verification method selected here.\\
NFR-02 & Core workflows support keyboard use, visible focus, meaningful control names, and a textual graph alternative. Target Web Content Accessibility Guidelines (WCAG) 2.2 level AA \citep{wcag}. Use manual and automated checks; automated scanning alone is insufficient.\par\smallskip \textbf{Sources: normative basis and project adaptation.} \evidence{E12}, WCAG 2.2 SC 1.1.1, 2.1.1, 2.4.7, 4.1.2 and Section 5.2.1 \citep{wcag}. The AA target, textual workflow and combined manual/automated checks are this project's adoption and test protocol.\\
NFR-03 & Proposed local selection, view-switching, and undo feedback target: 95th percentile of 200 ms within the display limits in FR-12 on the declared benchmark. Measure network requests separately.\par\smallskip \textbf{Sources: project performance target.} SemSpect, \evidence{E08}, Sections 1--2 \citep{semspect}, motivates complexity control. The percentile, 200-ms value, operations and display limits are selected engineering targets, not paper-derived thresholds.\\
NFR-04 & Proposed search, facet, and single-hop expansion target: 95th percentile of 2 s end-to-end on the controlled warm-source benchmark. Report cold runs, bounded paths, and public-endpoint variability separately. Arbitrary queries and external sources have no universal latency guarantee.\par\smallskip \textbf{Sources: documented motivation and project target.} Sparklis, \evidence{E04}, Section 5.4 \citep{sparklis}, explains endpoint-dependent response times. The 2-s percentile target and controlled warm-source conditions are project decisions.\\
NFR-05 & Default timeout is 10 s. Initial expansion retrieves at most 100 relationships; table pages contain at most 100 rows. Cancellation updates local state within 200 ms. Where source cancellation is unavailable, discard late results and report outstanding source execution accurately.\par\smallskip \textbf{Sources: documented precedent and project limits.} Sparklis, \evidence{E04}, Section 5.5 \citep{sparklis}, and LODmilla, \evidence{E05}, Section 3.2 \citep{lodmilla}, motivate bounded work. All stated timeout, count and cancellation values are selected here; late-result handling is a design extension.\\
NFR-06 & Keep credentials separate from workspace data. Never expose them in exports, URLs, histories, or logs. Stored credentials require a protected mechanism. Source permissions and operation controls enforce read-only access.\par\smallskip \textbf{Sources: security guidance and project policy.} \evidence{E13}, OWASP REST Security, \emph{Sensitive information in HTTP requests} \citep{owasprest}, recommends keeping authentication data out of URLs that may be captured in server logs. Protected storage, export/history exclusions and read-only operation controls are project security decisions under Section 1.\\
NFR-07 & Inject timeout, disconnect, malformed result, and expired authentication. The last valid workspace survives with query, source, constraints, completeness, and stale state.\par\smallskip \textbf{Sources: design extension and project test policy.} Tabulator, \evidence{E01}, Section 3.1 \citep{tabulator}, shows network failures; LODmilla, \evidence{E05}, Section 3.1 \citep{lodmilla}, saves graph state. The injected failure set and retained-state pass condition are specified here.\\
NFR-08 & Run mandatory tasks on at least two distinct domain vocabularies per backend without application-code changes. Presentation configuration may differ.\par\smallskip \textbf{Sources: documented motivation and project portability test.} Sparklis, \evidence{E04}, Section 5.5 \citep{sparklis}, discovers endpoint structure; /facet, \evidence{E02}, Sections 3.4 and 4.4 \citep{facet}, configures facets at runtime. Two vocabularies per backend and the no-code-change criterion are project decisions.\\
NFR-09 & Record fixture version, expected results, application and connector versions, dialect, language-reference release and commit, compatibility manifest, environment, cache conditions, failures, and timings for each acceptance run. Report backend results separately.\par\smallskip \textbf{Sources: reference-based verification and project reporting policy.} \evidence{E11}, the pinned openCypher TCK README \citep{opencypher}, describes scenario inputs and expected results/errors. The acceptance-run manifest, cache/timing metadata and separate backend reports extend that reference for this application.\\
NFR-10 & Document deployment configuration, allowed sources, adapter versions and supported browsers. Authenticated deployments use HTTPS except for isolated localhost testing. Query logging and evaluation recording are disabled by default; enabled logs redact authentication and source-sensitive values. Saving, exporting, importing and deleting local workspaces are available without enabling telemetry.\par\smallskip \textbf{Sources: project deployment and privacy policy.} \evidence{E13} \citep{owasprest} supports avoiding credential exposure; deployment and local-storage defaults are project decisions. Verify configuration, browser storage, representative request/log records and exported files in AC-10--11.\\
NFR-11 & At 320 CSS pixels wide and at 200\% browser zoom, core controls, source scope, error recovery and textual results remain operable. Wide native queries and graph canvases may have labelled local scrolling without forcing the whole page to scroll horizontally. Tested browser versions are declared in the acceptance manifest.\par\smallskip \textbf{Sources: normative basis and project viewport tests.} \evidence{E12}, WCAG 2.2 \citep{wcag}, supplies the adopted accessibility framework. The stated viewport and zoom cases are project checks supporting NFR-02; they do not alone establish AA conformance.\\
\end{longtable}}

\section{Acceptance and evaluation plan}
\subsection{Paired fixtures}
Provide RDF and property-graph representations of equivalent task facts, plus separate model-specific cases. Include ambiguous labels, disconnected entities, directed and inverse relationships, cycles, missing properties, empty strings, zero, false, temporal values, and known paths. RDF tests add typed and language-tagged literals, blank nodes, and named graphs. Property-graph tests add multiple labels, relationship properties, parallel relationships, and self-loops.

The paired fixtures, result oracles and pass conditions are project verification choices. Each scenario below lists the requirements it tests. Every task has expected identities, values, scope, and completeness. Compare task-level answers through a declared mapping without erasing model-specific distinctions. Equivalent tasks need not use identical query text or storage representation. Execute every scenario below on both connectors.

{\small
\begin{longtable}{@{}P{1.35cm}P{14.35cm}@{}}
\caption{Mandatory paired acceptance scenarios.}\label{tab:ac}\\
\toprule ID & Scenario and pass condition\\\midrule
\endfirsthead\toprule ID & Scenario and pass condition\\\midrule\endhead
\bottomrule\endlastfoot
AC-01 & Connect, choose scope, inspect equal-labelled entities and their incoming/outgoing edges, and attempt mutation. Identities and model-specific values are retained; entities remain separate; reads succeed and writes are rejected. Named graph and relationship-property cases pass independently.\par\smallskip \textbf{Project acceptance test derived from FR-01, FR-02, NFR-01 and NFR-06.}\par\smallskip \textbf{Sources.} Data identity and dataset scope: \evidence{E09}, RDF Concepts, Sections 3.1--3.4 and 4 \citep{rdfconcepts}; property-graph semantics: \evidence{E11}, the declared Cypher reference \citep{opencypher}. Read-only enforcement and mutation-rejection tests are project policy.\\
AC-02 & Search, disambiguate, combine categorical and range facets, remove a filter, and select missing values. Results match the oracle; OR/AND semantics, counts, scope restoration, completeness, and missing-value distinctions are correct.\par\smallskip \textbf{Project acceptance test derived from FR-03 and FR-04.}\par\smallskip \textbf{Sources.} Entity disambiguation: \evidence{E03}, RelFinder, Section 2.1 \citep{relfinder}; active facets and range filters: \evidence{E06}, Sampo-UI, Section 4.4 \citep{sampoui}. The oracle and pass conditions are project choices.\\
AC-03 & Expand, collapse, restore, and search a known path under direction/type/hop constraints. Neighbourhoods and paths match the oracle. Limits and omissions are visible; no-path and reached-limit outcomes remain distinct; undo preserves context.\par\smallskip \textbf{Project acceptance test derived from FR-05 and FR-06.}\par\smallskip \textbf{Sources.} Relationship paths and graph inspection: \evidence{E03}, RelFinder, Sections 2.2--2.3 \citep{relfinder}; state recovery and bounded traversal: \evidence{E05}, LODmilla, Sections 3.1--3.2 \citep{lodmilla}. Path limits and expected results are specified for this application.\\
AC-04 & Exercise every operator in Section~\ref{sec:profile} using the coverage cases below, independently on RDF/SPARQL and property-graph/Cypher. Inspect generated native text and edit a separate native query. Results and errors match the oracle; descriptions match semantics; native edits preserve visual state; unsupported translation is not fabricated. A compound example does not replace operator-level coverage.\par\smallskip \textbf{Project acceptance test derived from FR-07, FR-08 and NFR-01.}\par\smallskip \textbf{Sources.} Guided construction: \evidence{E04}, Sparklis, Sections 3 and 5.1 \citep{sparklis}; native-query inspection/editing: \evidence{E01}, Tabulator, Section 3.4 \citep{tabulator}. Operator semantics follow \evidence{E10}, SPARQL, Sections 5--9 and 11--17, and \evidence{E11}, the declared Cypher grammar/TCK reference \citep{sparql,opencypher}. The paired coverage cases are project tests, not a claim of complete TCK conformance.\\
AC-05 & Request and inspect guidance, accept or reject it, create an empty result, and remove a constraint. Suggestion certainty is explicit; chosen actions update query and description consistently; guidance remains optional; previous state survives.\par\smallskip \textbf{Project acceptance test derived from FR-09 and FR-10.}\par\smallskip \textbf{Sources.} Refinement and query verbalization: \evidence{E04}, Sparklis, Sections 3 and 5.2--5.3 \citep{sparklis}; ranked expansions: \evidence{E07}, Graph-Query Suggestions, Sections 3--4 \citep{suggestions}. Certainty labels and recovery pass conditions are application extensions.\\
AC-06 & Select an entity across views, filter, group results, and drill into an aggregate. Scope and selection remain consistent; counts reconcile with underlying exact or bounded data; textual controls support equivalent inspection.\par\smallskip \textbf{Project acceptance test derived from FR-11, FR-12 and NFR-02.}\par\smallskip \textbf{Sources.} Coordinated views: \evidence{E06}, Sampo-UI, Sections 4.4--4.5 \citep{sampoui}; aggregation: \evidence{E08}, SemSpect, Section 2 \citep{semspect}; accessible interaction: \evidence{E12}, WCAG SC 1.1.1, 2.1.1, 2.4.7 and 4.1.2 \citep{wcag}. Paired-result and selection checks are project tests.\\
AC-07 & Save/reopen a branch, replay after fixture change, and export results, query, view, and evidence bundle. State survives; change is flagged; identity, typed values, source, query, times, limits, and completeness remain correct; credentials are absent.\par\smallskip \textbf{Project acceptance test derived from FR-13, FR-14, FR-15 and NFR-06.}\par\smallskip \textbf{Sources.} Saved graph state: \evidence{E05}, LODmilla, Section 3.1 \citep{lodmilla}; source provenance: \evidence{E01}, Tabulator, Section 3.1 \citep{tabulator}; CSV export: \evidence{E06}, Sampo-UI, Section 4.1 \citep{sampoui}. Branching, change detection, additional export formats and credential exclusion are project extensions or policy.\\
AC-08 & Inject timeout, disconnect, authentication and decoding errors, and late cancelled results. Errors remain distinct; valid state survives with stale indication; retry succeeds; late results cannot overwrite state.\par\smallskip \textbf{Project acceptance test derived from FR-16, NFR-05 and NFR-07.}\par\smallskip \textbf{Sources.} Network feedback: \evidence{E01}, Tabulator, Section 3.1 \citep{tabulator}; endpoint-dependent performance and limits: \evidence{E04}, Sparklis, Sections 5.4--5.5 \citep{sparklis}. Failure injection, retry and late-response pass conditions are project tests.\\
AC-09 & Run display and performance workloads, exceed the display cap, and repeat on a second vocabulary. Grouping and limits work; declared benchmark targets are met; both backends and domains pass independently.\par\smallskip \textbf{Project acceptance test derived from FR-12, NFR-03, NFR-04 and NFR-08.}\par\smallskip \textbf{Sources.} Complexity control: \evidence{E08}, SemSpect, Sections 1--2 \citep{semspect}; endpoint conditions and portability: \evidence{E04}, Sparklis, Sections 5.4--5.5 \citep{sparklis}. Workloads, performance thresholds and repeated-domain pass conditions are project decisions.\\
AC-10 & Configure public and authenticated source profiles, reject unsupported required constructs, switch sources with an existing query, expire authentication, and inspect browser storage, requests, logs and exports. Source, graph scope and dialect remain visible; no old query runs on source change; secret material is absent from persisted or exported application data; an expired session preserves the workspace and offers reauthentication. Repeat core controls with keyboard, 320-pixel width and 200\% zoom.\par\smallskip \textbf{Project acceptance test derived from FR-01, FR-07, FR-18 and NFR-02, NFR-06, NFR-10--11.} Source and language fidelity follow \evidence{E09}--\evidence{E11}; accessibility follows \evidence{E12}; credential-exposure guidance is \evidence{E13}. Configuration, failure and viewport cases are project tests.\\
AC-11 & Export a workspace, clear local state, import into a fresh browser profile and compare every required field. Inject an unsupported schema version, wrong field types, an oversized file, markup in annotations and a mutating query. Invalid imports leave the active workspace intact; imported text is inert; execution still applies read-only controls; credentials are absent. Delete a saved workspace and verify associated local records are removed. A changed source requires an explicit new-workspace or replacement choice.\par\smallskip \textbf{Project acceptance test derived from FR-13--15, FR-19 and NFR-01, NFR-06--07, NFR-10.} Saved-state and provenance precedents are \evidence{E01} and \evidence{E05}; the interchange and failure-injection contract is a project decision.\\
\end{longtable}}

\subsection{Visual query coverage for AC-04}
Every operator below requires individual and applicable combined cases on both connectors. Fixtures declare expected identities, typed values, multiplicities, ordering, and completeness for each case. Include positive cases, empty-result cases, and validation cases for inapplicable value types. Each of the five aggregate functions is tested separately. Passing one aggregate or one backend does not satisfy the visual profile.

{\small
\begin{longtable}{@{}P{4.6cm}P{11.1cm}@{}}
\caption{Mandatory operator coverage within AC-04.}\label{tab:query-coverage}\\
\toprule Operators & Required checks on RDF/SPARQL and property-graph/Cypher\\\midrule
\endfirsthead\toprule Operators & Required checks on RDF/SPARQL and property-graph/Cypher\\\midrule\endhead
\bottomrule\endlastfoot
Connected patterns, variables, repeated-variable joins, type/label constraints & Build and edit each construct; verify shared-variable equality, join scope, type or label membership, and returned identities against the oracle. Equal labels do not substitute for equal entities.\\
Directed and inverse relationships & Query each direction separately and in joined patterns. Verify relation type, endpoints, and the expected effect of reversing direction.\\
Optional patterns and alternatives & Check a matched and unmatched optional branch, each alternative, and their combinations. Verify missing or unbound values, branch scope, and result multiplicity.\\
Equality, inequality, string matching, numeric ranges, temporal ranges & Test each operator separately and in combined filters, including range boundaries. Verify datatype-appropriate comparison and distinctions between missing values, empty strings, zero, and false.\\
Projection, distinct results, ordering, limits & Check selected variables, duplicates before and after distinct projection, ascending and descending ordering on declared values, and limits applied to the correct result scope.\\
Grouped \texttt{COUNT}, \texttt{SUM}, \texttt{AVG}, \texttt{MIN}, \texttt{MAX} & Test every function independently, then combine aggregation with joins, optional values, filters, projection, ordering, and limits. Check grouping keys, applicable datatypes, missing values, empty input, and multiplicity against model-specific expected results.\\
\end{longtable}}

\subsection{Benchmark protocol}
The project selects the fixture size, hardware, network bound, repetition count and benchmark conditions below. The fixture contains 100,000 identifiable entities and 1,000,000 relationship facts, plus task attributes. Report RDF triple counts and property-graph node, relationship, and property counts separately; different encodings do not imply equal storage cost.

The reference client has at least four processor cores and 16 GB memory. Record actual hardware, client runtime, source hardware/software, indexes, cache state, and network condition. Local timing excludes network waits. End-to-end query tests use a controlled source with documented round-trip time at most 100 ms. Collect at least 30 repetitions per declared operation after documented warm-up. Separate cold/warm distributions and retain failures.

Measure search, categorical/range facets, single-hop expansion, aggregation, bounded paths, view switching, undo, cancellation, and reopening. Report each backend independently. A workload failing on one backend does not satisfy the baseline.

\subsection{Task-based usability evaluation}
Use the same information needs for both user groups: resolve ambiguity, restrict a set, explain a path, construct a query, and preserve a finding. Define initial state, correct answers, stopping criteria, and unaided completion. Record completion, wrong turns, actions, time, and participants' explanation of results.

The project's formative evaluation targets at least 80\% unaided completion of predefined core tasks by non-experts. This is a target, not a study result. Report user groups and backends separately with participant experience. Settle recruitment, consent, and retention before collecting participant data. A small formative study does not establish population-wide effectiveness.

\section{Traceability and release criteria}
Use the individual origin statements for clause-level provenance. Table~\ref{tab:trace} summarises the mappings; its grouped references do not support every clause or demonstrate implementation or achieved targets.
{\small
\begin{longtable}{@{}P{2.35cm}P{3.6cm}P{5.05cm}P{4.3cm}@{}}
\caption{Requirements, sources, design considerations, and tests.}\label{tab:trace}\\
\toprule Requirements & Sources / consideration & Required response & Verification\\\midrule
\endfirsthead\toprule Requirements & Sources / consideration & Required response & Verification\\\midrule\endhead
\bottomrule\endlastfoot
FR-01--02 & \evidence{E04}, \evidence{E08}--\evidence{E11} / G06 & Dual read-only connections and preserved semantics. & AC-01,07,08; NFR-01,06,08.\\
FR-03--04 & \evidence{E02}, \evidence{E03}, \evidence{E04}, \evidence{E06} / G01 & Search, disambiguation, facets. & AC-02; NFR-01,02,04.\\
FR-05--06 & \evidence{E03}, \evidence{E05}, \evidence{E08} / G02 & Navigation, bounded paths, context. & AC-03,09; NFR-01,03--05.\\
FR-07--08 & \evidence{E01}, \evidence{E04}, \evidence{E10}--\evidence{E11} / G01,G06 & Visual profile and preserved native mode. & AC-04,08; NFR-01,06,08.\\
FR-09--10 & \evidence{E04}, \evidence{E07} / G03 & Grounded suggestions and readable state. & AC-05; NFR-01,02.\\
FR-11--12 & \evidence{E01}, \evidence{E06}, \evidence{E08}, \evidence{E12} / G02,G05 & Coordinated views and aggregation. & AC-06,09; NFR-01--05.\\
FR-13--15 & \evidence{E01}, \evidence{E04}--\evidence{E06}, \evidence{E08}--\evidence{E09}, \evidence{E11} / G04 & Saved context and faithful exports. & AC-07; NFR-01,06,07,09.\\
FR-16 & \evidence{E01}, \evidence{E04} / G02,G06 & Cancellation and recovery. & AC-08; NFR-05--07.\\
FR-12a,17 & \evidence{E01}, \evidence{E02}, \evidence{E04}, \evidence{E06} + project choices & Typed views and optional evaluation records. & Type/recording-control inspection; NFR-06,09.\\
FR-18--19 & \evidence{E01}, \evidence{E04}, \evidence{E05}, \evidence{E09}--\evidence{E13} + project choices & Inspectable configuration, versioned interchange and local deletion. & AC-10--11; NFR-01--02,06--07,10--11.\\
\end{longtable}}

Release requires every Must requirement and scenario to pass separately for RDF/SPARQL and property-graph/Cypher, including model-specific checks. A combined percentage shall not conceal a connector failure. Performance targets apply to declared conditions. An adopted target change updates the specification before release evaluation.

Future federation, additional views or query constructs, and collaboration must preserve the version 1 baseline. Architecture and stack choices must preserve the specified semantics and pass the behavioural tests.

\subsection{Acceptance evidence and release record}
For each Must requirement, the release record shall name its acceptance cases, fixture versions, responsible implementer and verifier, test outcome and reproducible evidence. Outcomes are pass, fail or not run; a missing result is not a pass. Functional and semantic cases require all expected results and errors to match; accessibility requires the manual and automated checks stated in NFR-02; performance requires operation-level distributions under the declared benchmark. No implementation results are asserted by this specification.

The release package shall include the source-profile template without secrets, dialect compatibility manifest, versioned workspace schema and examples, both domain fixtures in both graph models, fixture oracles, execution instructions, test results, benchmark records and a short user guide covering search, source scope, query inspection, bounded results, saved-state recovery and export. A verifier can reproduce each failed or passed case from those materials. Participant recruitment, consent and retention are settled before the optional usability evaluation; technical fixture tests do not substitute for that evaluation.

\begin{thebibliography}{99}
\setlength{\itemsep}{5pt}
\bibitem[Berners-Lee et al.(2006)]{tabulator}
Berners-Lee, T., Chen, Y., Chilton, L., Connolly, D., Dhanaraj, R., Hollenbach, J., Lerer, A., \& Sheets, D. (2006).
Tabulator: Exploring and analyzing linked data on the Semantic Web.
\emph{Proceedings of the 3rd International Semantic Web User Interaction Workshop}.
\url{https://www.cs.columbia.edu/~chilton/web/my_publications/BernersLeeTabulator2006.pdf}
\bibitem[Ferr\'e(2017)]{sparklis}
Ferr\'e, S. (2017). Sparklis: An expressive query builder for SPARQL endpoints with guidance in natural language.
\emph{Semantic Web, 8}(3), 405--418. \url{https://doi.org/10.3233/SW-150208}
\bibitem[Harris \& Seaborne(2013)]{sparql}
Harris, S., \& Seaborne, A. (Eds.). (2013). \emph{SPARQL 1.1 query language}. World Wide Web Consortium.
\url{https://www.w3.org/TR/2013/REC-sparql11-query-20130321/}
\bibitem[Heim et al.(2009)]{relfinder}
Heim, P., Hellmann, S., Lehmann, J., Lohmann, S., \& Stegemann, T. (2009).
RelFinder: Revealing relationships in RDF knowledge bases. In \emph{Semantic multimedia} (pp. 182--187). Springer.
\url{https://doi.org/10.1007/978-3-642-10543-2_21}
\bibitem[Hildebrand et al.(2006)]{facet}
Hildebrand, M., van Ossenbruggen, J., \& Hardman, L. (2006).
/facet: A browser for heterogeneous Semantic Web repositories. In \emph{The Semantic Web -- ISWC 2006} (pp. 272--285). Springer.
\url{https://doi.org/10.1007/11926078_20}
\bibitem[Ikkala et al.(2022)]{sampoui}
Ikkala, E., Hyv\"onen, E., Rantala, H., \& Koho, M. (2022).
Sampo-UI: A full stack JavaScript framework for developing semantic portal user interfaces.
\emph{Semantic Web, 13}(1), 69--84. \url{https://doi.org/10.3233/SW-210428}
\bibitem[Liebig et al.(2023)]{semspect}
Liebig, T., Opitz, M., Vialard, V., \& Wenzel, M. (2023).
Scalable no-code knowledge graph exploration and querying with SemSpect.
In \emph{Posters and Demo Track of the 19th International Conference on Semantic Systems} (CEUR Workshop Proceedings, Vol. 3526). CEUR-WS.org.
\url{https://ceur-ws.org/Vol-3526/paper-07.pdf}
\bibitem[Lissandrini et al.(2020)]{suggestions}
Lissandrini, M., Mottin, D., Palpanas, T., \& Velegrakis, Y. (2020).
Graph-query suggestions for knowledge graph exploration.
In \emph{Proceedings of The Web Conference 2020} (pp. 2549--2555). Association for Computing Machinery.
\url{https://doi.org/10.1145/3366423.3380005}
\bibitem[Micsik et al.(2014)]{lodmilla}
Micsik, A., T\'oth, Z., \& Turbucz, S. (2014).
LODmilla: Shared visualization of Linked Open Data.
In \emph{Theory and practice of digital libraries -- TPDL 2013 selected workshops} (pp. 89--100). Springer.
\url{https://doi.org/10.1007/978-3-319-08425-1_9}
\bibitem[openCypher(2026)]{opencypher}
openCypher. (2026). \emph{The openCypher property graph query language} (Release 2024.3) [Language specification and technology compatibility kit]. GitHub.
\url{https://github.com/opencypher/openCypher/releases/tag/2024.3}
\bibitem[OWASP Foundation(n.d.)]{owasprest}
OWASP Foundation. (n.d.). \emph{REST security cheat sheet}. OWASP Cheat Sheet Series. Retrieved October 6, 2026, from
\url{https://cheatsheetseries.owasp.org/cheatsheets/REST_Security_Cheat_Sheet.html}
\bibitem[World Wide Web Consortium(2014)]{rdfconcepts}
World Wide Web Consortium. (2014). \emph{RDF 1.1 concepts and abstract syntax}.
\url{https://www.w3.org/TR/2014/REC-rdf11-concepts-20140225/}
\bibitem[World Wide Web Consortium(2024)]{wcag}
World Wide Web Consortium. (2024). \emph{Web Content Accessibility Guidelines (WCAG) 2.2}.
\url{https://www.w3.org/TR/2024/REC-WCAG22-20241212/}
\end{thebibliography}

\appendix
\section{Local primary-source identifiers}
{\small
\begin{longtable}{@{}P{1.1cm}P{14.6cm}@{}}
\toprule ID & Local source\\\midrule
\endfirsthead\toprule ID & Local source\\\midrule\endhead
\bottomrule\endlastfoot
\evidence{E01} & \path{papers/2006_Berners-Lee_Tabulator_Exploring.pdf}\\
\evidence{E02} & \path{papers/2006_Hildebrand_facet_A.pdf}\\
\evidence{E03} & \path{papers/2009_Heim_RelFinder_Revealing.pdf}\\
\evidence{E04} & \path{papers/2017_Ferre_Sparklis_An.pdf}\\
\evidence{E05} & \path{papers/2014_Micsik_LODmilla_Shared.pdf}\\
\evidence{E06} & \path{papers/2022_Ikkala_Sampo-UI_A.pdf}\\
\evidence{E07} & \path{papers/2020_Lissandrini_Graph-Query_Suggestions.pdf}\\
\evidence{E08} & \path{papers/2023_Liebig_Scalable_No-Code.pdf}\\
\end{longtable}}
\section{Glossary}
\begin{description}[style=nextline,leftmargin=0pt,labelwidth=0pt,labelindent=0pt,itemindent=0pt]
\item[Entity] A source-identified object inspected in the interface.
\item[Relationship] An RDF predicate assertion or property-graph relationship with model-specific identity and context.
\item[Facet] A displayed dimension restricting the active result set.
\item[Workspace] Saved source references, query, constraints, display state, observations, and annotations.
\item[Exact result] A declared query scope evaluated without sampling or an undisclosed limit.
\item[Bounded result] Data restricted by an explicit count, depth, timeout, page, or sampling policy.
\item[Evidence bundle] A finding bound to source, query, parameters, retrieval time, scope, limits, and annotations.
\item[Paired acceptance] The same supported information need tested on both graph families, with separate model-specific checks.
\end{description}
\end{document}
